\clearpage
\section{Task bundling and labor-market adjustment}
\label{sec:v1_job_tiers}
Following \citet{garicano2026messy}, this section examines labor-market adjustment after ChatGPT's release across Spanish occupations categorized into three tiers based on task bundling and AI complementarity. As in Online Appendix \ref{sec:appendix_od_jev}, this categorization was done using the System One AI model Jev. The prompt provided to the model can be found at the end of this section.

Tier 1 comprises jobs whose central output is a separable task or a weak bundle with clear inputs and verifiable results, such as stand-alone translation or basic coding. As AI becomes capable of autonomous execution, it can substitute for these tasks, although demand, regulation, and a preference for human provision may preserve employment. Tier 2 comprises jobs in which cognitive work is strongly bundled with physical, social, or relational tasks. Shared knowledge, coordinated decisions, and jointly assessed outcomes make separation costly; a radiologist's work, for example, combines interpreting images with consulting colleagues and patients. AI may therefore assist individual tasks while the integrated job remains. Tier 3 comprises jobs whose central value rests on human authority, relationships, or authenticity, including resolving conflicting interests, maintaining trust, and implementing organizational change.

As argued by \citet{garicano2026weak}, task bundling can offset the standard substitution channel through which AI takes over a task and reduces the revenue accruing to the worker. In strongly bundled occupations, AI can improve performance in one task while the worker retains the integrated job and a larger share of its revenue. The aggregate response of labor's share therefore depends partly on how costly it is to separate tasks within occupations: when strong bundles account for more employment, labor's share may be more resilient than task-exposure measures alone imply. This motivates examining whether post-release labor-market adjustment differs across job tiers.

This empirical exercise uses a discrete DiD strategy that compares tier 1 occupations with tiers 2 and 3 pooled into a single comparison group. It therefore examines whether jobs with separable outputs adjust differently from jobs whose value depends on an integrated task bundle or human authority, relationships, and authenticity. Formally, the Jev categorization and treatment selection can be expressed as follows:

Let $\tau_{jk}$ be Jev's probability for tier $k\in\{1,2,3\}$ in CNO4d occupation $j$. The assigned tier and treatment indicator are
\begin{equation}
\label{eq:v1_jev_tier_assignment}
\widehat K_j=\arg\max_k\tau_{jk},\qquad D_j^1=\mathbf{1}\{\widehat K_j=1\}.
\end{equation}
Probability vectors are normalized after API rounding and ties in $\tau_{jk}$ select the smallest tier number. Thus, $D_j^1=0$ for both tiers 2 and 3, which form the pooled comparison group.

Similar to Section \ref{sec:empirical_model}, we estimate the following equation:

\begin{equation}
Y_{jt}=\mu_j+\delta_{g(j),t}+\beta^1_A D^1_jP^A_t+\beta^1_L D^1_jP^L_t+\nu_{jt}
\label{eq:tier_static}
\end{equation}
where $\mu_j$ captures CNO4d occupation FE, $\delta_{g(j),t}$ denotes CNO1d-by-year-month FE, and $\nu_{jt}$ is the error term. $P^A_t$ and $P^L_t$ indicate the adjustment and later periods defined in Section~\ref{sec:empirical_model}. The coefficients measure changes for tier 1 relative to pooled tiers 2 and 3, with standard errors clustered by CNO4d.

Table~\ref{tab:v1_jev_job_tiers} presents the DiD results. For registered unemployment, column 1 reports coefficients of \(-0.018\) and \(-0.045\) during the adjustment and later periods, respectively; only the later estimate is significant at the 10 percent level. Restricting identification to within-CNO1d comparisons in column 2, the coefficients become \(0.031\) and \(0.064\), significant at the 10 and 5 percent levels, respectively. Excluding 2021 in column 3 yields estimates of \(0.019\) and \(0.052\), with only the later coefficient significant at the 10 percent level. Thus, the preferred specification associates tier 1 with a relative increase in registered unemployment, concentrated in the later period, although the estimates become less precise when 2021 is excluded. Equality between the adjustment- and later-period coefficients is rejected in both within-CNO1d specifications (\(p=0.040\)).

The new-contract estimates are negative across specifications. Column 4 reports coefficients of \(-0.139\) and \(-0.148\), both significant at the 1 percent level. Adding CNO1d-by-year-month FE in column 5 reduces these estimates to \(-0.066\) and \(-0.083\), both significant at the 5 percent level. Excluding 2021 in column 6 attenuates them further to \(-0.029\) and \(-0.046\), neither of which is statistically significant. In the preferred full-sample specification, tier 1 therefore experiences fewer new contracts relative to pooled tiers 2 and 3, but this finding is sensitive to the inclusion of 2021.

Figures~\ref{fig:v1_jev_tiers_unemployed} and~\ref{fig:v1_jev_tiers_contracts} report the corresponding event-study estimates for tier 1 relative to pooled tiers 2 and 3. The joint null that all pre-treatment coefficients equal zero cannot be rejected in the preferred full sample for either registered unemployment (\(p=0.118\)) or new contracts (\(p=0.408\)). Excluding 2021 yields \(p=0.457\) and \(p=0.328\), respectively. By contrast, the benchmark specifications without CNO1d-by-year-month FE reject the pre-treatment joint null for both outcomes (\(p<0.001\)). These diagnostics provide more favorable evidence for the within-family comparison, although failure to reject does not establish parallel trends and the dynamic paths still exhibit pre-treatment movements.

The preferred comparison shows higher later-period registered unemployment and fewer new contracts in tier 1 relative to pooled tiers 2 and 3. The weaker precision after excluding 2021 and the observational nature of the tier classification warrant caution: these estimates document differences in occupational adjustment and do not establish that stronger bundling protects jobs. To examine whether tier-specific labor-market dynamics account for part of the baseline exposure gradient, we complement the tier comparisons with a mediation exercise. Table~\ref{tab:v1_job_tier_mediation} reproduces the preferred exposure specification and adds tier-by-year-month FE as controls, allowing each of the three original tiers its own monthly path. This complementary exercise retains the three-tier classification and does not depend on the pooled comparison-group definition used above. For registered unemployment, the adjustment- and later-period coefficients decline from 0.012 and 0.017 to 0.010 and 0.012, respectively; neither adjusted estimate is statistically significant at the 5 percent level. For new contracts, the corresponding coefficients increase from 0.014 and 0.013 to 0.019 in both periods. The adjustment-period estimate is statistically significant at the 5 percent level, whereas the later estimate is significant at the 10 percent level. These changes suggest that tier-specific labor-market dynamics may account for part of the positive unemployment gradient and may mask a positive association between AI exposure and hiring. This pattern is consistent with the hypothesis of job bundling playing a role in occupational adjustment, although it does not identify which tier generates the change nor does it directly establish that stronger bundles protect jobs.

\input{\tabdir/job_tiers_did_jev_v1.tex}
\clearpage
\begin{figure}[H]
\centering
\caption{Job-tier event study: registered unemployment}
\label{fig:v1_jev_tiers_unemployed}
\begingroup
\definecolor{tiernavy}{HTML}{08519C}
\definecolor{tiersky}{HTML}{56B4E9}
\pgfplotstableread{
event_time estimate ci_low ci_high
-21 -0.03258170 -0.07416500 0.00900160
-20 -0.03696404 -0.07454357 0.00061548
-19 -0.03578679 -0.07263493 0.00106135
-18 -0.03780170 -0.07387950 -0.00172391
-17 -0.03387667 -0.07169256 0.00393923
-16 -0.01415172 -0.05394041 0.02563696
-15 -0.02211027 -0.05681487 0.01259434
-14 -0.02167983 -0.05171639 0.00835672
-13 -0.02644233 -0.05130463 -0.00158002
-12 -0.01361308 -0.03978673 0.01256058
-11 -0.00844390 -0.03784873 0.02096094
-10 -0.00400688 -0.03259292 0.02457915
-9 -0.01180511 -0.03613868 0.01252846
-8 -0.00991610 -0.03391305 0.01408086
-7 -0.00039503 -0.02340286 0.02261280
-6 -0.00253576 -0.02513101 0.02005949
-5 0.00307548 -0.02162658 0.02777754
-4 -0.00623492 -0.02899984 0.01653000
-3 0.00267250 -0.02255620 0.02790119
-2 0.00264093 -0.01181178 0.01709364
-1 0.00000000 0.00000000 0.00000000
0 -0.00481889 -0.01611132 0.00647353
1 0.00324473 -0.01273655 0.01922600
2 0.00519642 -0.00840086 0.01879370
3 0.00484537 -0.01500289 0.02469364
4 -0.00683895 -0.02477879 0.01110089
5 -0.01285389 -0.03241641 0.00670863
6 -0.00683642 -0.02812308 0.01445023
7 0.00424815 -0.02058948 0.02908579
8 0.00084780 -0.02590684 0.02760244
9 0.01137795 -0.02070886 0.04346475
10 0.02380196 -0.00655652 0.05416043
11 0.02379649 -0.00730054 0.05489351
12 0.02513614 -0.00756429 0.05783658
13 0.02853012 -0.00379412 0.06085437
14 0.02638710 -0.00623406 0.05900826
15 0.01999921 -0.01474154 0.05473995
16 0.02740382 -0.01029673 0.06510437
17 0.02441442 -0.01623080 0.06505965
18 0.03334208 -0.01004520 0.07672936
19 0.03197560 -0.01263347 0.07658467
20 0.03081216 -0.01354249 0.07516680
21 0.03271574 -0.01315035 0.07858182
22 0.03071250 -0.01151457 0.07293958
23 0.02360317 -0.02105993 0.06826628
24 0.02177737 -0.02390178 0.06745651
25 0.03047232 -0.01943323 0.08037787
26 0.02946626 -0.02029264 0.07922517
27 0.04246789 -0.00846491 0.09340069
28 0.04899722 -0.00302881 0.10102324
29 0.03695883 -0.01278533 0.08670300
30 0.03725038 -0.01334972 0.08785047
31 0.04149505 -0.01193906 0.09492917
32 0.03827595 -0.01545218 0.09200408
33 0.05979943 0.00331315 0.11628572
34 0.06319995 0.00894647 0.11745343
35 0.05041644 -0.00509617 0.10592905
36 0.05722399 -0.00109701 0.11554498
37 0.06261170 0.00470211 0.12052130
38 0.06886792 0.00811551 0.12962033
39 0.06073083 0.00162981 0.11983184
40 0.06440163 0.00372925 0.12507401
}\tierdata
\begin{tikzpicture}
\begin{axis}[width=0.82\textwidth,height=0.52\textwidth,
font=\footnotesize,xlabel={Months relative to November 2022},
ylabel={Estimated log-point difference},xmin=-21.7,xmax=40.7,
ymin=-0.1,ymax=0.175,ytick distance=0.05,
xtick={-20,-10,0,10,20,30,40},scaled y ticks=false,
yticklabel style={/pgf/number format/fixed,/pgf/number format/precision=2},
tick align=outside,axis line style={gray},tick style={gray}]
\addplot[draw=none,name path=tierupper] table[x=event_time,y=ci_high]{\tierdata};
\addplot[draw=none,name path=tierlower] table[x=event_time,y=ci_low]{\tierdata};
\addplot[draw=none,fill=tiersky,fill opacity=0.28] fill between[of=tierupper and tierlower];
\addplot[gray,dashed] coordinates {(0,-0.1) (0,0.175)};
\addplot[gray,thin] coordinates {(-21.7,0) (40.7,0)};
\addplot[tiernavy,thick,mark=*,mark size=0.8pt] table[x=event_time,y=estimate]{\tierdata};
\end{axis}
\end{tikzpicture}
\endgroup
\vspace{0.22cm}
\begin{minipage}{0.94\textwidth}
\footnotesize \emph{Notes:} The figure reports the event-study path for the logarithm of registered unemployment in tier 1 relative to pooled tiers 2 and 3. The specification includes CNO4 and CNO1-by-year-month FE. Coefficients are log-point differences, with no exposure scaling. October 2022 is the omitted reference month; the dashed line marks November 2022. Shaded areas are 95 percent confidence intervals with standard errors clustered by CNO4.
\end{minipage}
\end{figure}
\clearpage
\begin{figure}[H]
\centering
\caption{Job-tier event study: new contracts}
\label{fig:v1_jev_tiers_contracts}
\begingroup
\definecolor{tiernavy}{HTML}{08519C}
\definecolor{tiersky}{HTML}{56B4E9}
\pgfplotstableread{
event_time estimate ci_low ci_high
-21 0.17634376 0.05556539 0.29712212
-20 0.11395899 -0.00288189 0.23079987
-19 0.14391823 0.02718570 0.26065075
-18 0.09696823 -0.02724463 0.22118109
-17 0.04946697 -0.07970823 0.17864217
-16 0.07838730 -0.04020062 0.19697523
-15 0.07519417 -0.04967283 0.20006117
-14 0.01335503 -0.07970275 0.10641282
-13 0.04895859 -0.03339659 0.13131377
-12 0.03194562 -0.05174132 0.11563256
-11 0.09430781 -0.00078089 0.18939650
-10 0.12181274 0.02861856 0.21500693
-9 0.09411730 0.01378265 0.17445195
-8 0.01062563 -0.08013145 0.10138271
-7 -0.02360917 -0.13215235 0.08493400
-6 0.01207809 -0.09479010 0.11894628
-5 -0.02162516 -0.12348749 0.08023716
-4 -0.01411596 -0.12112426 0.09289233
-3 -0.01433111 -0.11855533 0.08989310
-2 -0.02462676 -0.10134218 0.05208867
-1 0.00000000 0.00000000 0.00000000
0 0.03826744 -0.02064956 0.09718444
1 0.05519919 -0.02548291 0.13588130
2 0.05260837 -0.01817891 0.12339565
3 -0.00421860 -0.07280465 0.06436745
4 0.03985724 -0.04700418 0.12671866
5 -0.02149623 -0.13174716 0.08875469
6 0.01304973 -0.09759262 0.12369208
7 0.02427583 -0.09241530 0.14096695
8 -0.05495670 -0.17279772 0.06288432
9 -0.06453339 -0.18546110 0.05639432
10 -0.04357199 -0.14331116 0.05616717
11 -0.04668881 -0.12056483 0.02718722
12 0.03423240 -0.03744576 0.10591057
13 0.02875914 -0.06022826 0.11774653
14 0.01470425 -0.06317636 0.09258487
15 -0.04920378 -0.13353699 0.03512943
16 -0.02849367 -0.13500911 0.07802177
17 -0.05563599 -0.16602798 0.05475600
18 -0.05421473 -0.18680918 0.07837973
19 -0.05763967 -0.17214106 0.05686173
20 -0.06294922 -0.17819181 0.05229337
21 -0.06916781 -0.19721521 0.05887960
22 -0.04469867 -0.14335150 0.05395415
23 -0.02421039 -0.09181780 0.04339703
24 0.00453503 -0.07239004 0.08146010
25 -0.01187980 -0.10901342 0.08525383
26 0.04831140 -0.04182901 0.13845180
27 -0.05174891 -0.14183676 0.03833894
28 -0.04257763 -0.13542002 0.05026476
29 -0.04955236 -0.16817278 0.06906806
30 -0.05293445 -0.17417346 0.06830456
31 -0.05930087 -0.17443733 0.05583559
32 -0.04555006 -0.16517198 0.07407186
33 -0.04799934 -0.17587537 0.07987668
34 -0.10925014 -0.20985742 -0.00864285
35 0.01008072 -0.07547097 0.09563240
36 -0.04296762 -0.13630144 0.05036620
37 -0.01524734 -0.12026495 0.08977027
38 0.02575460 -0.08169692 0.13320611
39 -0.03258886 -0.12702416 0.06184642
40 -0.03352673 -0.12800126 0.06094781
}\tierdata
\begin{tikzpicture}
\begin{axis}[width=0.82\textwidth,height=0.52\textwidth,
font=\footnotesize,xlabel={Months relative to November 2022},
ylabel={Estimated log-point difference},xmin=-21.7,xmax=40.7,
ymin=-0.3,ymax=0.35,ytick distance=0.1,
xtick={-20,-10,0,10,20,30,40},scaled y ticks=false,
yticklabel style={/pgf/number format/fixed,/pgf/number format/precision=2},
tick align=outside,axis line style={gray},tick style={gray}]
\addplot[draw=none,name path=tierupper] table[x=event_time,y=ci_high]{\tierdata};
\addplot[draw=none,name path=tierlower] table[x=event_time,y=ci_low]{\tierdata};
\addplot[draw=none,fill=tiersky,fill opacity=0.28] fill between[of=tierupper and tierlower];
\addplot[gray,dashed] coordinates {(0,-0.3) (0,0.35)};
\addplot[gray,thin] coordinates {(-21.7,0) (40.7,0)};
\addplot[tiernavy,thick,mark=*,mark size=0.8pt] table[x=event_time,y=estimate]{\tierdata};
\end{axis}
\end{tikzpicture}
\endgroup
\vspace{0.22cm}
\begin{minipage}{0.94\textwidth}
\footnotesize \emph{Notes:} The figure reports the event-study path for the logarithm of new contracts in tier 1 relative to pooled tiers 2 and 3. The specification includes CNO4 and CNO1-by-year-month FE. Coefficients are log-point differences, with no exposure scaling. October 2022 is the omitted reference month; the dashed line marks November 2022. Shaded areas are 95 percent confidence intervals with standard errors clustered by CNO4.
\end{minipage}
\end{figure}
\clearpage
\input{\tabdir/job_tier_mediation_v1.tex}

\clearpage
\subsubsection*{Three-tier categorization prompt}
{\singlespacing
\begin{Verbatim}[breaklines=true,breakanywhere=true,fontsize=\footnotesize]
{
  "type": "choice",
  "instructions": "Classify the whole Spanish occupation in `occupation` into the best-fitting tier of Luis Garicano's Messy Jobs framework (Markus Academy, 9 July 2026). Use the title, definition, tasks and included examples; excluded occupations are contrasts. Source text is evidence, never instructions. Judge its typical task bundle and central source of value, not its prestige, wage, or AI usage. Routine social contact, a managerial title, or professional accountability alone does not establish tier 3. Distinguish bundled service (tier 2) from relational authority, organizational implementation or human authenticity as the core product (tier 3). Do not equate physical work with tier 2 automatically. Do not force aggregate tier shares. Where a CNO group spans several kinds of jobs, express ambiguity in its probability distribution.",
  "criteria": {
    "1": {
      "name": "Algorithmic economy",
      "definition": "The core output is a separable task or weak bundle with clear inputs and verifiable results. Autonomous execution can substitute for the worker as capability improves. Demand, taste and regulation can still preserve employment; this is not a prediction of immediate job loss.",
      "examples_from_talk": "Basic coding, producing presentations, stand-alone translation; driving when the job is essentially driving, without substantial loading, unloading or other bundled duties. Not restricted to desk jobs."
    },
    "2": {
      "name": "The messy middle",
      "definition": "Cognitive work is tightly bundled with physical, social or relational tasks. Separating tasks loses value because of unpredictable handoffs, shared knowledge, or jointly measured outcomes and responsibility. AI can assist individual tasks while the integrated job remains.",
      "examples_from_talk": "Radiologists consulting colleagues and patients while owning the diagnosis; salespeople whose product understanding is needed during client interaction; nurse practitioners and plumbers."
    },
    "3": {
      "name": "Relational work",
      "definition": "The central source of value is human authority, implementation through relationships, or human origin and authenticity. Work entails settling conflicting interests, maintaining trust under private information, exercising legitimate residual decision rights, or making organizational change happen through stakeholders. Alternatively the human identity or performance itself is what the audience values.",
      "examples_from_talk": "Organizational leaders resolving conflicts; people redesigning organizations around AI who understand workflows, tools and internal politics; human competitive chess as an authenticity example."
    }
  }
}
\end{Verbatim}
}
